\documentclass[11pt]{article}
\usepackage[a4paper,margin=2cm]{geometry}
\usepackage{xcolor}
\usepackage{enumitem}
\usepackage{booktabs}
\usepackage{tabularx}
\usepackage[most]{tcolorbox}
\usepackage{titlesec}
\usepackage{fancyhdr}
\usepackage{hyperref}
\usepackage{xurl}
\usepackage{needspace}
\newcolumntype{C}[1]{>{\centering\arraybackslash}p{#1}}

\definecolor{navy}{HTML}{002C5F}
\definecolor{sky}{HTML}{00AAD2}
\definecolor{sand}{HTML}{F3EFEA}
\definecolor{muted}{HTML}{6B7280}
\definecolor{ok}{HTML}{2E8B57}
\definecolor{alert}{HTML}{D9381E}
\definecolor{amber}{HTML}{C98A00}
\hypersetup{colorlinks=true,linkcolor=navy,urlcolor=navy,citecolor=navy}

\titleformat{\section}{\Large\bfseries\color{navy}}{\thesection.}{6pt}{}[{\color{sky}\titlerule[1.5pt]}]
\titleformat{\subsection}{\large\bfseries\color{navy}}{}{0pt}{}
\setlength{\parindent}{0pt}
\setlength{\parskip}{4pt}
\setlist{itemsep=1pt,topsep=2pt}
\pagestyle{fancy}\fancyhf{}
\fancyhead[L]{\small\color{muted}Honk-Sense: problem validity, use cases and research}
\fancyhead[R]{\small\color{muted}Piyush Anand \textbullet\ Ankith Deegoju}
\fancyfoot[C]{\small\thepage}
\renewcommand{\headrulewidth}{0pt}

\newtcolorbox{keybox}[1][]{enhanced,breakable,colback=sand,colframe=sand,boxrule=0pt,left=8pt,right=8pt,top=5pt,bottom=5pt,borderline west={3pt}{0pt}{sky},#1}
\newcommand{\strong}{{\color{ok}\bfseries Strong}}
\newcommand{\medium}{{\color{amber}\bfseries Medium}}
\newcommand{\weak}{{\color{alert}\bfseries Supporting only}}
\newcommand{\IN}{{\color{sky}\bfseries[India]}}
\newcommand{\chk}{\textsuperscript{\color{alert}\dag}}
\renewcommand{\cite}[1]{{\color{navy}[#1]}}

\begin{document}

\begin{tcolorbox}[colback=navy,colframe=navy,coltext=white,left=10pt,right=10pt,top=8pt,bottom=8pt]
{\LARGE\bfseries Honk-Sense}\\[2pt]
{\large Problem validity, use cases and research findings}\\[6pt]
\small HCoE Mobility Innovation Ideathon 2026 \textbullet\ Topic 1: India-Smart Interaction \textbullet\ Track A\\
Piyush Anand \textbullet\ Ankith Deegoju \textbullet\ IIT Hyderabad
\end{tcolorbox}

\begin{keybox}
\textbf{The short answer to our central question.}
Mirrors, dashcams, parking sensors and radar blind-spot monitors all \emph{look} for objects. None of them \emph{listen}. On Indian roads the horn and the siren are how other road users tell you their intent: ``I am here'', ``I am overtaking'', ``give way''. A sealed, air-conditioned cabin with music on blocks much of that sound, and even when a driver hears it, they often cannot tell where it came from.

Honk-Sense adds value in four places where existing systems cannot help:
\begin{enumerate}[leftmargin=*]
  \item \textbf{Emergency vehicles}: the direction and approach of an ambulance before it is visible.
  \item \textbf{Sound from out of sight}: a horn around a blind turn or behind a bus or truck. A siren carries where cameras and radar cannot see.
  \item \textbf{Direction of a horn}: which side a honking two-wheeler or truck is on, in cars that have no radar blind-spot monitoring, which is most cars sold in India today.
  \item \textbf{Accessibility}: drivers with hearing loss, who have been eligible for licences in India since 2016.
\end{enumerate}

Where existing systems are already enough, Honk-Sense should \emph{not} claim value: parking, silent vehicles in the blind spot (radar's job), and recording evidence (dashcam's job). Our honest position is that Honk-Sense \textbf{complements} the existing sensors and does not replace any of them.
\end{keybox}

{\small\color{muted}\textbf{How to read this document.} \IN\ marks studies and data from India. \chk\ marks a number we found through secondary sources (news or review articles). Check it against the primary report before putting it on a slide. Bracketed codes such as \cite{R1} point to the reference list in Section~12.}

% =====================================================
\section{Objective of the discussion: is the problem valid?}
We asked whether the Honk-Sense problem is real and whether it is relevant to Hyundai. The evidence says \textbf{yes, with a narrow and honest scope}:
\begin{itemize}[leftmargin=*]
  \item \textbf{Cabins block outside warnings.} A closed car body cuts 21--35\,dB at 1\,kHz, the band where horns and sirens sit \cite{R12}\chk. In a US DOT study, drivers with windows closed found the right direction of a siren only \textbf{26\%} of the time, and sirens were effective only \textbf{25--40\,ft (8--12\,m)} ahead at intersections \cite{R11}\chk.
  \item \textbf{Music makes it worse.} For every 2\,dB of in-car music, the level a siren needs to be heard rises by 1.0--1.7\,dB \cite{R13}.
  \item \textbf{Drivers miss what they should see.} 75\% of lane-change and merge crashes involve a \emph{recognition failure}: the driver looked but did not register the other vehicle \cite{R15}.
  \item \textbf{India uses sound as a signal more than most countries.} ``Horn OK Please'' is painted on trucks, ``Blow Horn'' signs mark hill turns, and two-wheelers honk to announce themselves. Two IIT/Microsoft India papers built entire traffic-sensing systems on honks for exactly this reason \cite{R6,R7}\,\IN.
  \item \textbf{Delays to ambulances cost lives.} The Law Commission estimated that about \textbf{50\%} of road deaths could be avoided with care in the first hour \cite{R4}\chk\,\IN. A Mumbai survey of ambulance drivers reported that about 65\% of private cars do not make way \cite{R5}\chk\,\IN.
\end{itemize}
The problem is \emph{not} ``drivers cannot see cars''; cameras and radar already address that. The problem is: \textbf{important sounds outside the car do not reach the driver in a usable form}, meaning on time, with a direction, and separated from the noise.

% =====================================================
\section{Existing systems and their limitations}
\renewcommand{\arraystretch}{1.25}
{\small
\begin{tabularx}{\textwidth}{@{}p{2.6cm}XXX@{}}
\toprule
\textbf{System} & \textbf{What it solves} & \textbf{Limitations} & \textbf{Where Honk-Sense adds value} \\
\midrule
Rear-view and side mirrors & Cheap, always there; shows the area behind and to the side & Blind spots remain. Only works if the driver looks: in naturalistic data, 46\% looked at the left mirror before a left lane change, and most skipped the over-the-shoulder check \cite{R16,R17}. Poor at night and in rain & A sound cue that does not need the driver to be looking in the right place \\
Dashcam & Records evidence for insurance and disputes & Records only; gives no live warning. Usually faces forward only & None as a recorder. Honk-Sense is a live warning, not a recording \\
Ultrasonic parking sensors & Low-speed obstacles while parking & Range of about 2\,m, and active only at low speed \cite{R20} & Not relevant for parking; ultrasonic is enough there \\
Radar blind-spot monitoring (BSM) & Vehicles in the adjacent lane, including silent ones. Cuts lane-change crashes by about 14\% \cite{R14} & Available mostly on upper variants in India. Struggles in dense mixed traffic, with two-wheelers weaving and faded lane markings \cite{R18,R19}\,\IN. Blind to vehicles hidden behind other vehicles & Adds intent (a horn means ``I am here'' or ``I am coming through''). Hears what is behind an obstruction. Lower cost for cars without BSM \\
Camera ADAS (lane departure, forward collision) & Lanes, vehicles ahead, pedestrians & Needs line of sight. Harder in rain, glare and dust & Hears an ambulance before any camera can see it \\
Sirens and horns themselves & Warn other road users & Weakened by the cabin, masked by music, and hard to place by direction \cite{R11,R12,R13} & Exactly the gap: brings the warning inside, with its direction \\
\bottomrule
\end{tabularx}}

\textbf{What still goes wrong despite all of these.}
Every system above is either \emph{visual} (it needs line of sight) or \emph{passive} (it needs the driver to look). None of them carries the \emph{intent} that a horn or siren expresses, and none of them brings that intent into the cabin with a direction.

% =====================================================
\section{Proposed use cases}
We kept the ideas from our discussion, made them sharper, and rated each one honestly. A rating of \strong\ means existing systems clearly cannot do this. \medium\ means Honk-Sense helps, but another system partly covers it. \weak\ means it is a secondary benefit, not a headline.

\subsection{U1. Emergency-vehicle direction and approach \hfill \strong}
\textbf{Situation.} An ambulance is 100--200\,m behind, hidden by a bus. The driver has music on and windows closed.\\
\textbf{Why existing systems fail.} Cameras and radar cannot see it. The siren reaches the cabin late, and drivers place it correctly only about a quarter of the time \cite{R11}.\\
\textbf{What Honk-Sense does.} Detects the siren pattern, shows ``Ambulance, behind left, approaching'', lowers the music, and suggests which way to move.\\
\textbf{Honest limit.} The detection range depends on traffic noise, so we should not promise ``200\,m''. Promise \emph{earlier than the driver would notice}, and validate the actual distance in tests.

\subsection{U2. Direction of a horn: who is honking, and where? \hfill \strong}
\textbf{Situation.} City traffic, horns from every side. One of them is a bike right at your rear-left quarter as you start to turn left.\\
\textbf{Why existing systems fail.} The driver hears a horn but cannot tell which one matters or where it is. After hearing hundreds of horns, drivers stop reacting to any of them.\\
\textbf{What Honk-Sense does.} Shows only the horns that are \emph{close and directed at you}, with an arrow, and ignores the rest. This filtering is the feature.

\subsection{U3. Blind-spot awareness when the other vehicle honks \hfill \medium}
\textbf{Situation.} A two-wheeler or truck in your blind spot honks before you change lanes.\\
\textbf{Why existing systems fail.} Mirrors miss the blind spot, and most cars in India have no radar BSM.\\
\textbf{What Honk-Sense does.} ``Horn, right rear'' just as you signal right.\\
\textbf{Honest limit.} It only works if the other vehicle honks. Radar BSM catches \emph{silent} vehicles too, so say ``complements BSM'', not ``replaces BSM''. In Indian traffic, though, vehicles overtaking or filtering usually do honk.

\subsection{U4. Out of sight: blind turns, ghat roads, behind large vehicles \hfill \strong}
\textbf{Situation.} A hairpin bend with a ``Blow Horn'' sign. A bus is coming the other way and honks before it is visible.\\
\textbf{Why existing systems fail.} Every visual sensor needs line of sight; sound travels around the corner.\\
\textbf{What Honk-Sense does.} ``Vehicle ahead, around the bend'', several seconds before it is visible.

\subsection{U5. Drivers with hearing loss \hfill \strong\ (for this group)}
\textbf{Situation.} India has allowed hearing-impaired people to get driving licences since a MoRTH circular in 2016, with further guidance in 2023 \cite{R8}\IN.\\
\textbf{Why existing systems fail.} Horns and sirens are audio-only warnings.\\
\textbf{What Honk-Sense does.} Turns them into visual cues on the cluster and mirror, and optionally a haptic cue in the seat or steering wheel. This is an inclusion story judges will remember.

\subsection{U6. Fewer interruptions: filtering the noise \hfill \medium}
\textbf{Situation.} Peak-hour traffic in Hyderabad, where noise often exceeds 90--100\,dB \cite{R9,R10}\IN.\\
\textbf{Value.} Less stress and less ``alarm fatigue''. Honk-Sense stays silent for honks not meant for you and speaks up only for the ones that matter.

\subsection{U7. Level crossings and rural roads \hfill \weak}
Train horns and pressure horns of trucks at night. This is useful, but Indian Railways has removed most unmanned level crossings on its main network, so keep it as a supporting example only.

\subsection{U8. Fleet and city data \hfill \weak}
Anonymous counts of where vehicles honk and where sirens get stuck can feed city planning. Nericell \IN\ showed this was feasible with phones in 2008 \cite{R6}. Treat it as a later-stage idea; it raises privacy questions.

\textbf{Recommended top five for the pitch: U1, U2, U4, U5 and U3}, in that order. U1 and U4 are the strongest because no visual sensor can do them at all.

% =====================================================
\needspace{14\baselineskip}\section{Comparison with existing technologies}
{\small\renewcommand{\arraystretch}{1.25}
\setlength{\tabcolsep}{4pt}\begin{tabularx}{\textwidth}{@{}XC{1.4cm}C{1.6cm}C{1.85cm}C{1.4cm}C{1.4cm}C{1.5cm}@{}}
\toprule
\textbf{Situation} & \textbf{Mirrors} & \textbf{Dashcam} & \textbf{Ultrasonic} & \textbf{Radar BSM} & \textbf{Camera ADAS} & \textbf{Honk-Sense} \\
\midrule
Ambulance hidden behind traffic & -- & -- & -- & -- & -- & \textbf{Yes} \\
Horn from behind a blind turn & -- & -- & -- & -- & -- & \textbf{Yes} \\
Which side a horn came from & -- & -- & -- & -- & -- & \textbf{Yes} \\
Honking bike in the blind spot & Partly & -- & -- & Yes & -- & \textbf{Yes} \\
\emph{Silent} car in the blind spot & Partly & -- & -- & \textbf{Yes} & -- & No \\
Parking and low-speed obstacles & Partly & -- & \textbf{Yes} & -- & Partly & No \\
Lane departure, vehicle ahead & -- & -- & -- & -- & \textbf{Yes} & No \\
Evidence after a crash & -- & \textbf{Yes} & -- & -- & -- & No \\
Driver with hearing loss & Yes & -- & -- & Yes & Yes & \textbf{Yes} \\
\bottomrule
\end{tabularx}}

\textbf{Where Honk-Sense gives a bigger contribution:} anything that is \emph{heard before it is seen}, and anything where the \emph{direction of a sound} matters.\\
\textbf{Where existing systems are enough:} silent vehicles, parking, lane keeping and evidence. Saying this openly in Q\&A builds trust.

% =====================================================
\section{Research papers and existing studies}

\subsection{Acoustic sensing of honks: Indian work}
\begin{itemize}[leftmargin=*]
  \item \textbf{Nericell} (Mohan, Padmanabhan, Ramjee; Microsoft Research India; ACM SenSys 2008) \IN. Phones detected honks and braking on Bengaluru roads, with audio processed on the phone for privacy. It won the ACM Test-of-Time award in 2019 \cite{R6}. \emph{Use:} proves that honk detection in Indian traffic is feasible and that on-device processing is the accepted privacy approach.
  \item \textbf{Horn-OK-Please} (Sen, Raman, Sharma; IIT Bombay; ACM MobiSys 2010) \IN. Estimated vehicle speed from the Doppler shift of honks, using over 18 hours of roadside recordings \cite{R7}. \emph{Use:} shows that honks carry useful motion information, not just presence.
\end{itemize}

\subsection{Siren detection and direction finding}
\begin{itemize}[leftmargin=*]
  \item \textbf{Listening for Sirens} (Marchegiani and Newman, Oxford, 2018): classifies sirens and horns in city noise and estimates their direction with two microphones \cite{R21}.
  \item \textbf{Emergency Vehicles Audio Detection and Localization in Autonomous Driving} (Sun et al., 2021): uses low-cost microphones on a car to estimate siren presence, direction and distance in real driving \cite{R22}.
  \item \textbf{SirenNet} (Tran and Tsai, IEEE Access 2020): separates sirens, horns and noise with 98\% accuracy on its test set \cite{R23}.
  \item \textbf{Acoustic Localization of Emergency Vehicle Sirens} (Ohlenbusch et al., DAGA 2020): an eight-microphone array tested on real traffic \cite{R24}.
  \item \textbf{Fraunhofer IDMT ``Hearing Car''}, with CARIAD (VW Group): external microphone modules at the rear of the car, AI detection and siren tracking. Tested over 1,500\,km in Sweden in 2025 \cite{R25}. \emph{Use:} shows the industry is moving here, while no Indian-traffic product exists.
  \item \textbf{Interior versus exterior microphones} for horns and sirens (Forum Acusticum 2025) \cite{R26}: relevant to where we mount the microphones.
  \item \textbf{Hybrid framework for emergency vehicles on Indian roads} (IEEE 2020, Bengaluru) \IN: camera plus siren pitch detection for traffic-signal priority \cite{R27}. It is infrastructure-side, not in-car, so it leaves our space open.
  \item \textbf{GCC-PHAT} (Knapp and Carter, 1976): the standard method for finding direction from the time difference between microphones.
\end{itemize}

\subsection{How audible sirens and horns are inside a car}
\begin{itemize}[leftmargin=*]
  \item US DOT audible-warning study (1970s): effective siren range at an urban intersection of 25--40\,ft; drivers with windows closed found the correct direction 26\% of the time \cite{R11}\chk.
  \item US National Bureau of Standards, 1974: a closed car body removes 21--35\,dB at 1\,kHz \cite{R12}\chk.
  \item Baba and Ebata (Acoustical Society of America, 1996): every 2\,dB of music raises the siren detection threshold by 1.0--1.7\,dB \cite{R13}.
  \item Cerence driver study (2019): 67\% of drivers wanted to be told immediately when a siren is detected and wanted the audio turned down \cite{R28}. This is a survey of preferences, but it supports our ``duck the music'' feature.
\end{itemize}

\subsection{Blind spots and lane changes}
\begin{itemize}[leftmargin=*]
  \item Cicchino (IIHS), \emph{Traffic Injury Prevention} 2018: blind-spot monitoring is linked to \textbf{14\%} fewer lane-change crashes and 23\% fewer with injuries; the injury figure is less statistically certain \cite{R14}.
  \item FMCSA: \textbf{75\%} of lane-change and merge crashes involve a recognition failure by the driver changing lanes \cite{R15}.
  \item NHTSA 100-Car naturalistic study (DOT HS 811 147, 2009): before a planned left lane change, only 46\% of drivers glanced at the left mirror \cite{R16}.
  \item Side blind-zone alert study: drivers skipped the over-the-shoulder glance in 68\% of left and 85\% of right lane changes. With alerts on, lane changes without a mirror check fell by 31\% \cite{R17}.
\end{itemize}

\subsection{ADAS in Indian conditions}
\begin{itemize}[leftmargin=*]
  \item Kumbhar and Pyasi, SAE 2026-26-0035 \IN: field test of lane-departure warning, AEB and blind-spot monitoring on Indian roads. Effectiveness drops with poor lane markings and unpredictable traffic \cite{R18}.
  \item SAE 2026-26-0042 review and ARAI commentary \IN: unmarked lanes, dense two-wheeler traffic and false alerts are the main challenges for ADAS in India \cite{R19}.
\end{itemize}

{\small\color{muted}\textbf{Gap we found:} no published in-car system detects horns and sirens with direction, trained on Indian traffic sound. Indian work so far is either roadside traffic sensing (Nericell, Horn-OK-Please) or signal-priority infrastructure. This is the clearest evidence that Honk-Sense is new.}

% =====================================================
\needspace{14\baselineskip}\section{Indian road conditions and traffic behaviour}
{\small\renewcommand{\arraystretch}{1.25}
\begin{tabularx}{\textwidth}{@{}p{3.4cm}XX@{}}
\toprule
\textbf{Condition} & \textbf{What is different in India} & \textbf{Opportunity for Honk-Sense} \\
\midrule
Mixed vehicles & Two-wheelers are about 45\% of road deaths \cite{R1}\IN. Autos, buses, trucks, cycles and carts share one road & Two-wheelers and autos weave close and honk; direction matters most here \\
Horn as language & ``Horn OK Please'' on trucks; ``Blow Horn'' signs on hills \cite{R7}\IN & The horn carries intent. Honk-Sense reads that intent \\
Very loud streets & Traffic noise of 90--100\,dB or more measured in Mumbai \cite{R9,R10}\IN & Drivers tune out all horns, so filtering only the relevant ones is valuable \\
Weak lane discipline & Faded or missing markings; vehicles filter between lanes & Camera and radar ADAS lose accuracy \cite{R18}\IN; sound does not depend on lane markings \\
Ghat and hill roads & Hairpin bends with no line of sight & Horn heard around the bend, before the vehicle is visible (U4) \\
Congested cities & Ambulances stuck in traffic; drivers do not know which way to move \cite{R5}\IN & Direction and a suggested move (U1) \\
Urban and rural highways & Trucks with pressure horns; night driving; music on long trips & Long-trip family drivers with music on are the most cut off \\
Monsoon and fog & Rain on the roof, glare, low visibility & Visual sensors degrade; sound still travels (noise filtering must handle rain) \\
\bottomrule
\end{tabularx}}

% =====================================================
\needspace{14\baselineskip}\section{Accident statistics and supporting data}
{\small\renewcommand{\arraystretch}{1.25}
\begin{tabularx}{\textwidth}{@{}X>{\raggedright\arraybackslash}p{4.4cm}>{\raggedright\arraybackslash}p{4.0cm}@{}}
\toprule
\textbf{Data point} & \textbf{Number} & \textbf{Source} \\
\midrule
Road accidents in India, 2023 & 4,80,583 accidents; 1,72,890 deaths (about 20 deaths an hour) & MoRTH, \emph{Road Accidents in India 2023} \cite{R1}\IN \\
Two-wheeler riders' share of deaths, 2023 & 44.8\% (about 77,500 deaths) & MoRTH 2023 \cite{R1}\IN \\
Pedestrians' share of deaths, 2023 & 20.4\% & MoRTH 2023 \cite{R1}\IN \\
Hit from back (rear-end) & About 21\% of accidents, 20\% of deaths & MoRTH 2023 \cite{R2}\chk\IN \\
Hit from side & About 15\% of accidents & MoRTH 2023 \cite{R2}\chk\IN \\
Head-on collisions & About 17\% of accidents & MoRTH 2023 \cite{R2}\chk\IN \\
Over-speeding, 2022 & 72.3\% of accidents; 71.2\% of deaths & MoRTH 2022 \cite{R3}\chk\IN \\
Driving on the wrong side, 2022 & About 5\% of accidents & MoRTH 2022 \cite{R3}\chk\IN \\
Deaths avoidable with first-hour care & About 50\% & Law Commission, 201st Report \cite{R4}\chk\IN \\
Private cars not giving way to ambulances & About 65\% (Mumbai, survey of over 100 ambulance drivers) & Radhee Foundation, 2018 \cite{R5}\chk\IN \\
Lane-change crashes with recognition failure & 75\% & FMCSA \cite{R15} \\
Effect of blind-spot monitoring & 14\% fewer lane-change crashes & IIHS \cite{R14} \\
\bottomrule
\end{tabularx}}

\textbf{What the data does \emph{not} show, and how to say it.}
\begin{itemize}[leftmargin=*]
  \item MoRTH has \textbf{no ``blind spot'' or ``failed to hear'' category}. Causes are recorded as speeding, wrong side and so on. Say so openly: \emph{``India does not record these causes yet; side and rear-end collisions together make up over a third of crashes, and that is where awareness of the surroundings matters.''}
  \item There is \textbf{no national figure for ambulance delays caused by traffic}. Use the Law Commission's golden-hour estimate and the Mumbai survey, labelled as one city.
  \item \textbf{Do not claim a percentage reduction for Honk-Sense.} The closest analogy is blind-spot monitoring (14\%) and the 31\% drop in unchecked lane changes when alerts were on \cite{R14,R17}. Say: ``Awareness systems have shown 14--31\% improvements in similar situations; we would measure ours in a pilot.''
\end{itemize}

% =====================================================
\section{Driver psychology and behaviour}
\begin{itemize}[leftmargin=*]
  \item \textbf{Habituation and alarm fatigue.} When everything honks, the brain stops responding to horns, including the important ones. Too many alerts cause the same problem with car systems, and drivers switch them off. ARAI studies how false alerts erode trust \cite{R19}. \emph{Design rule:} Honk-Sense must be silent most of the time.
  \item \textbf{Looked but failed to see.} Drivers can look straight at a mirror and still miss a bike; 75\% of lane-change crashes involve this kind of recognition failure \cite{R15}. A sound cue with a direction works even when visual attention is elsewhere.
  \item \textbf{Locating sounds inside a car.} People locate sounds well in the open, but a cabin blurs direction, especially front versus back. Only 26\% got the siren direction right \cite{R11}. The result is freezing or moving the wrong way when an ambulance comes.
  \item \textbf{Masking by music and conversation.} Music raises the level a siren needs to be heard \cite{R13}. Family trips, the main use case for Hyundai cars, have the most cabin noise.
  \item \textbf{Mirror-checking habits.} Most drivers skip the over-the-shoulder glance \cite{R16,R17}. Two-wheelers filtering between lanes make this more dangerous in India.
  \item \textbf{Stress in congestion.} Constant honking raises stress and irritability, and stressed drivers make more errors. Removing irrelevant honks from the driver's attention is a comfort and safety benefit (U6).
  \item \textbf{Reaction to warnings.} Drivers respond faster to a warning that tells them \emph{what} and \emph{where} than to a general beep. This is why Honk-Sense shows a direction arrow and a label (``Ambulance, behind left''), not just a chime.
\end{itemize}

% =====================================================
\needspace{14\baselineskip}\section{Installation possibilities}
{\small\renewcommand{\arraystretch}{1.25}
\begin{tabularx}{\textwidth}{@{}>{\raggedright\arraybackslash}p{3.2cm}>{\raggedright\arraybackslash}X>{\raggedright\arraybackslash}X@{}}
\toprule
\textbf{Location} & \textbf{Advantages} & \textbf{Problems} \\
\midrule
Side-mirror housings (left and right) & Wide spacing between left and right, so good direction. Already has wiring and is fairly sheltered & Wind noise at speed; needs a wind cover \\
Rear bumper or tailgate & Low wind noise (Fraunhofer's choice \cite{R25}); good for sounds from behind (ambulances, overtaking) & Road splash and dirt; needs sealing \\
Front bumper or grille & Hears oncoming horns and sounds at blind turns & Engine and fan noise; stone chips \\
Roof (shark-fin antenna) & 360\textdegree\ view; out of splash & Microphones close together, so direction is less precise; wind noise \\
Inside the cabin & Cheapest; no weather sealing & The cabin blocks and colours outside sound; cabin speech and music get mixed in \cite{R26} \\
\midrule
\multicolumn{3}{@{}l}{\textbf{Where the driver sees and hears the alert}} \\
\midrule
Instrument cluster or HUD & Arrow and label in the line of sight & Must not add clutter \\
Side-mirror LED & Natural place to look before a lane change & Left and right only \\
Directional cabin audio & Plays a short cue from the speaker on the correct side; lowers music & Must be brief, not another alarm \\
Seat or steering vibration & Works for drivers with hearing loss (U5) & Optional; adds cost \\
\bottomrule
\end{tabularx}}

\textbf{Recommendation:} four microphones, one at each corner (two in the mirror housings and two at the rear corners), plus an optional front one. This gives direction around the whole car while keeping wind noise manageable. Alerts go to the cluster and mirror LEDs, with a short directional sound cue. This matches the ``4 microphones'' already in our deck.

% =====================================================
\section{Why accidents still happen despite existing systems}
\begin{enumerate}[leftmargin=*]
  \item \textbf{Mirrors depend on attention.} Drivers do not check them, or check and fail to register \cite{R15,R16}.
  \item \textbf{Blind-spot systems are not universal.} Radar BSM is mostly on upper variants, and in India it is tuned for orderly lanes, not weaving two-wheelers \cite{R18,R19}.
  \item \textbf{Every sensor needs line of sight.} Nothing sees around a bus, a truck or a hairpin bend.
  \item \textbf{Proximity sensors are short-range and low-speed.} Ultrasonic sensors work up to about 2\,m and are made for parking \cite{R20}.
  \item \textbf{Sirens lose to the cabin.} Sealing, AC and music cut them down, and drivers cannot tell where they come from \cite{R11,R12,R13}.
  \item \textbf{Too many horns means no horn.} Habituation makes the meaningful honk invisible.
  \item \textbf{The gap no one addresses:} \emph{turning outside sound into a short, directional, filtered cue inside the car.} That is Honk-Sense.
\end{enumerate}

% =====================================================
\needspace{14\baselineskip}\section{Expected research outcome and next steps}
\begin{tabularx}{\textwidth}{@{}Xl@{}}
\toprule
\textbf{Outcome we wanted} & \textbf{Status} \\
\midrule
Validity of the problem statement & Done (Sections 1 and 10) \\
At least five valuable use cases & Done: 8 use cases, top 5 chosen (Section 3) \\
Comparison with mirrors, dashcams, sensors and BSM & Done (Sections 2 and 4) \\
Research papers, preferably Indian & Done: 2 core Indian papers plus international work (Section 5) \\
Indian accident statistics & Done, with gaps stated honestly (Section 7) \\
Driver psychology & Done (Section 8) \\
Installation locations & Done, with a recommendation (Section 9) \\
\bottomrule
\end{tabularx}

\textbf{To make it stronger before the final round:}
\begin{enumerate}[leftmargin=*]
  \item Check every \chk\ number against the primary report; the MoRTH PDF is the most important.
  \item Run a short driver survey of 30--50 people at IIT Hyderabad and in our families. Ask: How often do you hear a horn and not know where it is? When did you last notice an ambulance late? Would you want the music to duck for a siren? This replaces ``based on our experience'' with real data.
  \item Record about 30 minutes of traffic sound from a phone on the dashboard (inside) and outside the window, in Hyderabad traffic and on a highway. Even a simple spectrogram showing how much the cabin blocks sound makes a strong slide.
\end{enumerate}

% =====================================================
\section{References}
{\small
\begin{description}[leftmargin=1.2cm,style=nextline,itemsep=2pt]
\item[R1] MoRTH, \emph{Road Accidents in India 2023} (as reported by PTI). \url{https://www.dtnext.in/news/national/172-lakh-people-died-in-480-lakh-road-accidents-in-2023-road-ministry-report-844831}
\item[R2] Collision-type shares, 2023 (from MoRTH via PTI). \url{https://theprint.in/economy/not-wearing-seat-belts-claimed-16025-lives-in-road-accidents-in-2023-morth/2731093}
\item[R3] MoRTH, \emph{Road Accidents in India 2022}: over-speeding and wrong-side shares. \url{https://www.lokmattimes.com/national/overspeeding-accounted-for-712-of-total-deaths-in-road-accidents-morth/amp/}
\item[R4] Law Commission of India, 201st Report (golden hour), as cited by Deccan Herald. \url{https://www.deccanherald.com/archives/auto-drivers-may-become-life-2063220}
\item[R5] Radhee Foundation survey of ambulance drivers, Mumbai, 2018 (DNA). \url{https://www.dnaindia.com/mumbai/report-motorists-on-mumbai-s-roads-are-insensitive-ambulance-drivers-2614716}
\item[R6] Mohan, Padmanabhan, Ramjee. \emph{Nericell: Rich monitoring of road and traffic conditions using mobile smartphones}. ACM SenSys 2008. \url{https://www.microsoft.com/en-us/research/project/nericell/}
\item[R7] Sen, Raman, Sharma. \emph{Horn-Ok-Please}. ACM MobiSys 2010, DOI 10.1145/1814433.1814449. \url{https://api.crossref.org/works/10.1145\%2F1814433.1814449}
\item[R8] Driving licences for hearing-impaired persons in India (MoRTH 2016 circular; Delhi HC). \url{https://hrln.org/litigation/deaf-eligible-to-undertake-driving-test-licence-delhi-hc}
\item[R9] The levels and effects of traffic noise: a study in Mumbai (IIT Bombay). \url{https://www.audiology-worldnews.com/awareness/1844-the-levels-and-effects-of-traffic-noise-a-study-in-mumbai-india}
\item[R10] MPCB sub-group minutes on noise standards (horn and siren levels). \url{https://mpcb.gov.in/sites/default/files/internal-study-groups/Minutes_2ndMeetingSub\%20Group_noise\%20standards_siren.pdf}
\item[R11] US DOT siren audibility study, as summarised in JEMS. \url{https://www.jems.com/ems-operations/ground-ambulance-operations/why-motorists-dont-get-out-way-0/}
\item[R12] NBS 1974 car-body sound reduction, as summarised in Fire Engineering. \url{https://www.fireengineering.com/fire-apparatus/apparatus-siren-limitations-and-intersection-crashes/}
\item[R13] Baba and Ebata, effect of in-car music on siren detection, ASA 1996. \url{https://auditory.org/asamtgs/asa96haw/5aPP/5aPP15.html}
\item[R14] Cicchino, \emph{Effects of blind spot monitoring systems on police-reported lane-change crashes}, Traffic Injury Prevention 2018. \url{https://iihs.org/topics/bibliography/ref/2143}
\item[R15] FMCSA, inadequate surveillance and lane-change crashes. \url{https://webarchive.library.unt.edu/web/20161101172143/https://cms.fmcsa.dot.gov/safety/driver-safety/cmv-driving-tips-inadequate-surveillance}
\item[R16] NHTSA, \emph{Lane change crashes in the 100-Car study}, DOT HS 811 147, 2009. \url{https://www.nhtsa.gov/sites/nhtsa.gov/files/811147.pdf}
\item[R17] \emph{Lane change behavior with a side blind zone alert system} (SWOV listing). \url{https://swov.nl/en/publicatie/lane-change-behavior-side-blind-zone-alert-system}
\item[R18] Kumbhar and Pyasi, \emph{Actual user experience and efficacy of ADAS features in Indian driving conditions}, SAE 2026-26-0035. \url{https://trid.trb.org/View/2761739}
\item[R19] ARAI on ADAS validation for Indian roads. \url{https://www.motorindiaonline.in/how-arai-is-building-indias-adas-backbone-through-data-driven-validation-simulation-engg/}
\item[R20] Texas Instruments, ultrasonic versus radar sensing (SPRY331). \url{https://www.ti.com/lit/pdf/spry331}
\item[R21] Marchegiani and Newman, \emph{Listening for Sirens}, 2018. \url{https://arxiv.org/abs/1810.04989}
\item[R22] Sun et al., \emph{Emergency Vehicles Audio Detection and Localization in Autonomous Driving}, 2021. \url{https://arxiv.org/abs/2109.14797}
\item[R23] Tran and Tsai, \emph{SirenNet}, IEEE Access 2020. \url{https://doaj.org/article/5080316787ee43de91bdb588704ad98f}
\item[R24] Ohlenbusch et al., \emph{Acoustic localization of emergency vehicle sirens in traffic}, DAGA 2020. \url{https://www.pub.dega-akustik.de/DAGA_2020/data/articles/000010.pdf}
\item[R25] Fraunhofer IDMT, \emph{The Hearing Car} (IEEE Spectrum). \url{https://spectrum.ieee.org/fraunhofer-hearing-car-microphone}
\item[R26] Detection of horns and sirens with interior versus exterior microphones, Forum Acusticum 2025. \url{https://dael.euracoustics.org/confs/fa2025/data/articles/000006.pdf}
\item[R27] \emph{A hybrid framework for expediting emergency vehicle movement on Indian roads}, IEEE 2020. \url{https://ieeexplore.ieee.org/document/9074933}
\item[R28] Cerence, emergency vehicle detection driver study, 2019. \url{https://cerence.gcs-web.com/news-releases/news-release-details/hear-call-help-how-emergency-vehicle-detection-can-make-your-car}
\end{description}}

\end{document}
